\chapter{Die Model-Landschaft 2026 -- Ein Vergleich}
\label{chap:model_landschaft}

% ============================================================
% LERNZIELE
% ============================================================
\begin{lernziele}
Nach diesem Kapitel kannst du:
\begin{itemize}
    \item Die drei führenden Closed-Source-Anbieter (OpenAI, Anthropic, Google) mit ihren Stärken und Schwächen vergleichen
    \item Die wichtigsten Open-Source-Modelle (Llama, Mistral, Qwen, DeepSeek) nach Kriterien wie Leistung, Lizenz und Betriebskosten bewerten
    \item Eine fundierte Modell-Entscheidung basierend auf Anforderungen wie Qualität, Kosten, Latenz und Datenhoheit treffen
    \item Einen Multi-Provider-Client implementieren, der den Austausch zwischen Modellen abstrahiert
    \item Vendor-Lock-in-Risiken identifizieren und eine Ausweichstrategie entwerfen
\end{itemize}
\end{lernziele}

% ============================================================
% MOTIVATION
% ============================================================
\section{Motivation}

2026 gibt es nicht \glqq das eine beste Modell\grqq{}. Es gibt Dutzende Modelle, die sich in Preis, Leistung, Latenz, Kontextfenster und Sicherheit stark unterscheiden. Die Wahl des falschen Modells kostet Unternehmen täglich tausende Euro -- entweder durch überhöhte API-Kosten (zu teures Modell), durch schlechte Ergebnisse (falsche Modell-Familie) oder durch Ausfallrisiken (Single-Provider-Abhängigkeit).

Ein AI Engineer muss die Landschaft nicht nur kennen, sondern aktiv steuern können: Wann setzt man GPT-4o-mini ein, wann Llama 3 lokal, wann DeepSeek R1 für Reasoning? Welche Modelle tauscht man aus, wenn der Preis steigt? Welcher Provider bietet EU-Datenhaltung?

Dieses Kapitel gibt dir einen strukturierten Überblick über den Modell-Markt 2026, eine Entscheidungsmatrix für die tägliche Arbeit und die Werkzeuge, um Modell-Entscheidungen systematisch zu treffen.

\section{Die Architektur des modernen AI-Stacks}

Bevor wir die einzelnen Modelle vergleichen, lohnt ein Blick auf das Gesamtsystem. Ein LLM allein ist selten die Lösung -- es ist eine Komponente in einem größeren Stack. 2026 hat sich eine vierschichtige Architektur als Standard etabliert:

\begin{verbatim}
   +------------------------------------------------------------+
   |              Die vier Schichten des AI-Stacks              |
   |                                                            |
   |  1. LLM (Thinker)        Reasoning, Wissen, Kreativitaet   |
   |       |                                                    |
   |  2. RAG (Knowledge)      Kontext, Doku, Datenbanken        |
   |       |                                                    |
   |  3. Agents (Workers)     Ausfuehrung, Automatisierung      |
   |       |                                                    |
   |  4. MCP (Infrastructure) Verbindung zu Tools, APIs, FS     |
   +------------------------------------------------------------+
\end{verbatim}

\subsection{Layer 1: Das LLM als Denker}

Das LLM ist der Kern. Es kann argumentieren, schreiben, analysieren und Probleme lösen -- aber es kennt nur das, worauf es trainiert wurde. Es hat keinen Zugriff auf Unternehmensdaten, Live-Informationen oder externe Tools. Es besitzt reines Wissen und Reasoning.

Die Wahl des richtigen LLMs (die wir im Rest dieses Kapitels vertiefen) bestimmt die fundamentale Qualität des gesamten Stacks. Ein schwaches Basismodell limitiert alle darüber liegenden Schichten.

\subsection{Layer 2: Retrieval-Augmented Generation}

RAG gibt dem LLM Zugriff auf eine Wissensbasis -- Unternehmensdokumente, Richtlinien, Datenbanken, interne Wikis, aktuelle Kundeninformationen. Vor jeder Antwort wird relevantes Wissen abgefragt und als Kontext in den Prompt eingefügt.

\begin{verbatim}
   Ohne RAG:      [Frage] --> [LLM (nur Trainingswissen)] --> [Antwort]
   
   Mit RAG:       [Frage] --> [Vector Search] --> [Relevante Doku]
                                |
                                v
                   [LLM (Prompt + Doku)] --> [Antwort mit Fakten]
\end{verbatim}

RAG ist das Fundament für faktentreue Antworten. Kapitel~\ref{chap:rag_vector_dbs} behandelt RAG-Systeme ausführlich.

\subsection{Layer 3: Agents als Arbeiter}

Ein Agent (Kapitel~\ref{chap:ai_agents}) hört nicht bei Textgenerierung auf. Er führt Aktionen aus: Tickets lesen, E-Mails senden, Berichte erstellen, CRMs aktualisieren, Code deployen, Workflows steuern. Agents verwandeln Intelligenz in Ausführung.

Der entscheidende Unterschied zu einem reinen RAG-System: Agents haben einen Loop (Observe -- Plan -- Act -- Reflect) und können sich über mehrere Schritte erstrecken.

\subsection{Layer 4: MCP als Infrastruktur}

Der \textbf{Model Context Protocol} (MCP) ist die Schicht, die die meisten Teams übersehen. MCP ist der Standard, der KI-Systeme mit Tools, Datenbanken, APIs, Dateien und externen Diensten verbindet. Man kann es als USB-C für KI betrachten: ein Protokoll, tausende Integrationen, kein Custom-Connector für jedes Tool.

MCP standardisiert:

\begin{itemize}[leftmargin=*]
    \item \textbf{Tool Discovery}: Der Client erfragt verfügbare Tools vom MCP-Server
    \item \textbf{Tool Execution}: Der Client ruft ein Tool mit standardisierten Parametern auf
    \item \textbf{Resource Access}: Lese-/Schreibzugriff auf Dateien, Datenbanken, APIs
    \item \textbf{Prompts}: Vorgefertigte Prompt-Vorlagen auf dem Server
\end{itemize}

\begin{verbatim}
   Ohne MCP:        [LLM] --[Custom Code]--> [API A]
                      | --[Custom Code]--> [DB B]
                      | --[Custom Code]--> [Tool C]
                      
   Mit MCP:          [LLM] --[MCP Client]--> [MCP Server] --> [API A]
                                               |           --> [DB B]
                                               |           --> [Tool C]
\end{verbatim}

\subsection{Zusammenspiel der Schichten}

Die eigentliche Revolution ist nicht der bessere Chatbot -- es ist die Kopplung von Intelligenz mit Wissen, Werkzeugen und Arbeitsabläufen:

\begin{verbatim}
   Einfacher Chat:         LLM
   
   Chat mit Kontext:       LLM + RAG
   
   Automatisierung:        LLM + RAG + Agents
   
   Produktionssystem:      LLM + RAG + Agents + MCP
\end{verbatim}

Jede zusätzliche Schicht erhöht die Leistungsfähigkeit, aber auch die Komplexität, Latenz und Fehleranfälligkeit. Die Kunst liegt darin, nur die Schichten zu aktivieren, die für den konkreten Anwendungsfall notwendig sind.

Die folgenden Abschnitte dieses Kapitels helfen dir, die richtige \textbf{erste Schicht} -- das LLM -- für deinen Stack zu wählen. Die übrigen Schichten werden in den Kapiteln~\ref{chap:rag_vector_dbs},~\ref{chap:ai_agents} und im MCP-Exkurs vertieft.

% ============================================================
% GRUNDLAGEN
% ============================================================
\section{Grundlagen -- Modell-Klassifikation verstehen}

Bevor wir einzelne Modelle vergleichen, müssen wir die Metriken verstehen, anhand derer Modelle bewertet werden.

\subsection{Parameter -- die Modellgröße}

Die Anzahl der Parameter eines Modells ist die grobe Maßzahl für seine Kapazität. Ein Parameter ist ein Gewicht im neuronalen Netz. Mehr Parameter bedeuten mehr Speicherplatz für gelerntes Wissen, aber auch höhere Rechenkosten:

\begin{table}[H]
\centering
\begin{tabularx}{\linewidth}{lXX}
\toprule
\textbf{Größenklasse} & \textbf{Parameter} & \textbf{Typisches Einsatzgebiet} \\
\midrule
Klein & 1--8 Mrd. & Edge-Geräte, Mobile, Echtzeit, einfache Klassifikation \\
Mittel & 8--70 Mrd. & Allrounder, Production-Systeme, selbst gehostet \\
Groß & 70--400 Mrd. & Hochleistungs-Inference, Cloud-APIs \\
Sehr groß (MoE) & 400+ Mrd. & Spitzenmodelle, nur API verfügbar \\
\bottomrule
\end{tabularx}
\caption{Modellgrößenklassen und ihre Einsatzgebiete}
\label{tab:modellgroessen}
\end{table}

\subsection{MoE -- Mixture of Experts}

Ein Architektur-Muster, bei dem das Modell aus vielen \glqq Experten\grqq{} (Sub-Netzwerken) besteht. Für jede Anfrage werden nur die relevanten Experten aktiviert, nicht das gesamte Netzwerk. Das Ergebnis: Die Leistung eines großen Modells bei den Kosten eines kleineren.

\begin{verbatim}
  Ohne MoE:        Mit MoE:
  [LLM]            [Router] --+-- [Expert A] (Mathe)
                   (Anfrage)  +-- [Expert B] (Code)
                               +-- [Expert C] (Text)
                               +-- [Expert D] (Bild)

  Jede Anfrage:    Jede Anfrage:
  Alle 70B aktiv   Nur 7B von 70B aktiv
  -> hohe Kosten   -> niedrige Kosten
\end{verbatim}

\textbf{Beispiele:} GPT-4o (8 Experten, 220B total, 22B aktiv), Llama 4 Maverick (1T total, 16B aktiv).

\subsection{Benchmarks -- mit Vorsicht zu genießen}

Benchmarks wie MMLU, HumanEval, GSM8K oder Chatbot Arena geben einen Hinweis auf die Modellqualität, sind aber kein Ersatz für eigene Tests:

\begin{itemize}[leftmargin=*]
    \item \textbf{MMLU}: 57 Fächer von Medizin bis Jura -- misst allgemeines Wissen
    \item \textbf{HumanEval}: Code-Generierung aus Docstrings -- misst Programmierfähigkeit
    \item \textbf{GSM8K}: Mathematische Textaufgaben -- misst logisches Reasoning
    \item \textbf{Chatbot Arena}: Elo-Rating basierend auf menschlichen Präferenzen -- misst subjektive Qualität
\end{itemize}

\textbf{Problem:} Modelle werden auf Benchmarks optimiert (Benchmark-Overfitting). Ein hoher MMLU-Score garantiert nicht, dass ein Modell in deiner spezifischen Domäne gut ist.

\subsection{Weitere Schlüsselbegriffe}

\begin{description}[leftmargin=*,style=nextline]
    \item[Quantisierung] Reduktion der Parameter-Genauigkeit von FP16 auf INT8 oder FP4. Reduziert Speicherbedarf um 50--75\,\% bei geringem Qualitätsverlust. Unverzichtbar für Self-Hosting.
    \item[Context Window] Die maximale Token-Anzahl, die ein Modell in einem Durchlauf verarbeiten kann. 2026 zwischen 8K (klein) und 2M (Experimental).
    \item[Knowledge Cutoff] Das Datum, bis zu dem das Modell trainiert wurde. Modelle mit Cutoff 2025 kennen keine Ereignisse von 2026.
    \item[Open Weight] Die Modellgewichte sind öffentlich verfügbar, aber die Trainingsdaten sind es nicht.
    \item[Open Source] Modellgewichte \textit{und} Trainingscode/-daten sind öffentlich. Sehr selten bei großen Modellen.
\end{description}

% ============================================================
% DER MARKT IM ÜBERBLICK (existiert)
% ============================================================
\section{Der Markt im Überblick}

Im Jahr 2026 hat sich die Model-Landschaft von einem offenen zu einem oligopolistischen Feld entwickelt, in dem drei Akteure (\textbf{OpenAI}, \textbf{Anthropic} und \textbf{Google}) die Closed-Source-Spitze dominieren, während \textbf{Meta (Llama)}, \textbf{Mistral} und \textbf{Qwen} den Open-Source-Markt steuern.

\begin{table}[H]
\centering
\begin{tabularx}{\linewidth}{lXr}
\toprule
\textbf{Anbieter} & \textbf{Flaggschiff} & \textbf{Marktposition} \\
\midrule
OpenAI & GPT-4o, o-series & Marktführer, breitestes Ökosystem \\
Anthropic & Claude Opus 4 & Beste Safety, längster Kontext bei Sonnet \\
Google & Gemini 2.5 Ultra & Beste Multimodalität, Mega-Kontext \\
Meta & Llama 4 Maverick & Open-Source-Marktführer \\
Mistral & Mistral Large 2 & Bester EU-Anbieter, DSGVO-konform \\
Alibaba & Qwen 2.5 & Stärkste Konkurrenz aus Asien \\
DeepSeek & R1 / V3 & Preisbrecher, Reasoning auf o1-Niveau \\
\bottomrule
\end{tabularx}
\caption{Marktübersicht der wichtigsten Model-Anbieter 2026}
\label{tab:markt}
\end{table}

% ============================================================
% CLOSED SOURCE (existiert)
% ============================================================
\section{Closed-Source: Die Großen Drei}

\subsection{OpenAI -- GPT-4o und die o-Series}

OpenAI ist der Marktführer mit dem weitesten API-Ökosystem. Ihre Modelle lassen sich in drei Kategorien einteilen:

\begin{itemize}[leftmargin=*]
    \item \textbf{GPT-4o (Omni):} Multimodales Flaggschiff (Text, Audio, Bild), stark in der allgemeinen Kompetenz und Geschwindigkeit. Der Standard für die meisten Use-Cases.
    \item \textbf{o-Series (o1, o3):} Modelle mit erweiterten Reasoning-Fähigkeiten durch systematisches Grübeln vor der Antwort. Gut für mathematische, wissenschaftliche und komplexe logische Aufgaben -- aber langsamer und teurer.
    \item \textbf{GPT-4o-mini:} Der kleine Bruder -- 80\,\% der Qualität von GPT-4o zu 15\,\% der Kosten. Die erste Wahl für Production-Workloads.
\end{itemize}

\subsection{Anthropic -- Claude-Familie}

Anthropics Claude-Reihe ist bekannt für beste Safety und längere Kontextfenster:

\begin{itemize}[leftmargin=*]
    \item \textbf{Claude Opus 4:} Anthropics stärkstes Modell. Besonders gut in komplexer Analyse, Code-Generierung und kreativer Textproduktion.
    \item \textbf{Claude Sonnet 4:} Der Sweet Spot -- ausgezeichnete Qualität zu wettbewerbsfähigem Preis. Die Standardwahl für die meisten Produktivsysteme.
    \item \textbf{Claude Haiku:} Anthropics schnellstes und günstigstes Modell. Für einfache Klassifizierung und Filterung -- dort, wo Geschwindigkeit wichtiger ist als Nuance.
\end{itemize}

Besondere Stärken von Claude: \textbf{Native Tool Use} und das \textbf{200K Context Window}.

\subsection{Google -- Gemini-Familie}

Googles Gemini-Reihe punktet mit Multimodalität und der tiefsten Cloud-Integration (GCP):

\begin{itemize}[leftmargin=*]
    \item \textbf{Gemini 2.5 Ultra:} Googles Flaggschiff, konkurrenzfähig mit GPT-4o in Benchmarks. Besonders stark in Multimodalität (Text+Bild+Audio+Video).
    \item \textbf{Gemini Pro:} Das Arbeitstier -- gut genug für 90\,\% der Production-Einsätze, zu einem Bruchteil des Preises.
    \item \textbf{Gemini Flash:} Googles schnellstes Modell. Ideal für Echtzeit-Anwendungen.
    \item \textbf{Gemini Nano:} On-Device-Modell für mobile und Edge-Einsatzfälle.
\end{itemize}

Besonderheit: Gemini unterstützt Kontextfenster bis zu 1 Million Tokens -- das längste aller kommerziellen Modelle.

% ============================================================
% ARCHITEKTUR
% ============================================================
\section{Architektur -- Multi-Provider-Routing}

Ein produktives System verlässt sich nie auf einen einzigen Modell-Provider. Die folgende Architektur zeigt ein Gateway, das Anfragen je nach Anforderung an verschiedene Modelle routet:

% ABBILDUNG: Multi-Provider-Routing
\begin{verbatim}
                  +-----------------------+
                  |      User Request       |
                  +----------+------------+
                             |
                    +--------v--------+
                    | Request Classifier|
                    | (einfach/schwer/   |
                    |  kreativ/Code)    |
                    +--------+--------+
                             |
           +-----------------+------------------+
           |                 |                  |
      (einfach)        (schwer/Reasoning)   (kreativ/Code)
           |                 |                  |
     +-----v-----+    +-----v-----+      +-----v-----+
     | GPT-4o-    |    | Claude    |      | Codestral |
     | mini       |    | Opus 4    |      |           |
     +-----------+    +-----------+      +-----------+
           |                 |                  |
           +-----------------+------------------+
                             |
                    +--------v--------+
                    | Response         |
                    | Validator        |
                    | + Cost Logger    |
                    +--------+--------+
                             |
                    +--------v--------+
                    | Fallback: wenn   |
                    | ein Provider     |
                    | ausfallt         |
                    +-----------------+
\end{verbatim}

\subsection{Komponenten des Routings}

\begin{itemize}[leftmargin=*]
    \item \textbf{Request Classifier}: Bestimmt anhand der Anfrage, welches Modell am besten geeignet ist. Einfache Fragen $\rightarrow$ günstiges Modell. Komplexe Aufgaben $\rightarrow$ teures Modell.
    \item \textbf{Load Balancer}: Verhindert Rate-Limits durch Verteilung auf mehrere Provider.
    \item \textbf{Cost Logger}: Erfasst die Kosten pro Anfrage und Modell. Grundlage für Optimierung.
    \item \textbf{Fallback Chain}: Fällt ein Provider aus, wird automatisch der nächste kontaktiert. Ohne Unterbrechung für den User.
\end{itemize}

% ============================================================
% THEORIE
% ============================================================
\section{Theorie -- Warum es nicht ein bestes Modell gibt}

\subsection{Die Ökonomie der Modell-Anbieter}

Jeder Modell-Anbieter verfolgt eine andere Strategie:

\begin{itemize}[leftmargin=*]
    \item \textbf{OpenAI}: Plattform-Strategie. Verdient am API-Volumen. Ziel: möglichst viele Entwickler an die Plattform binden.
    \item \textbf{Anthropic}: Safety-first. Positioniert sich als Premium-Anbieter für Unternehmen, die Sicherheit und Kontrolle brauchen.
    \item \textbf{Google}: Integration in GCP. Verdient an Cloud-Infrastruktur, nicht primär am Modell.
    \item \textbf{Meta}: Open-Source als Hebel. Verdient am Llama-Ökosystem (Hardware-Partner, Cloud-Deals). Ziel: den Markt für KI-Infrastruktur öffnen.
    \item \textbf{Mistral}: EU-Nische. Setzt auf DSGVO-Konformität und europäische Sprachqualität als Differenzierungsmerkmal.
    \item \textbf{DeepSeek}: Preis-Krieg. Unterbietet alle etablierten Anbieter massiv, um Marktanteile zu gewinnen.
\end{itemize}

\subsection{Trade-offs bei der Modell-Wahl}

\begin{table}[H]
\centering
\begin{tabularx}{\linewidth}{lX}
\toprule
\textbf{Kriterium} & \textbf{Trade-off} \\
\midrule
Qualität vs. Kosten & Teurere Modelle liefern bessere Ergebnisse. Der Grenznutzen sinkt aber rapide. \\
Latenz vs. Qualität & Reasoning-Modelle (o1, DeepSeek R1) brauchen Sekunden bis Minuten für die Antwort. \\
Datenhoheit vs. Komfort & Self-Hosting bedeutet Betriebskosten und Wartung, aber volle Kontrolle. \\
Provider-Stabilität vs. Feature-Geschwindigkeit & Große Anbieter sind stabiler, kleine innovieren schneller. \\
API vs. SDK vs. Self-Hosted & APIs sind am einfachsten, SDKs bieten mehr Kontrolle, Self-Hosting maximiert Flexibilität. \\
\bottomrule
\end{tabularx}
\caption{Trade-offs bei der Modell- und Provider-Wahl}
\label{tab:provider-tradeoffs}
\end{table}

\subsection{Wann wechselt man den Provider?}

\begin{itemize}[leftmargin=*]
    \item \textbf{Preiserhöhung} des aktuellen Providers
    \item \textbf{Verschlechterung} der Modellqualität (heimliches Downgrade hinter generischem Model-String)
    \item \textbf{Compliance-Anforderungen} (neue Regulierung erfordert EU-Datenhaltung)
    \item \textbf{Ausfall} oder Rate-Limits des aktuellen Providers
    \item \textbf{Innovation} eines Konkurrenten (z.\,B. DeepSeek R1 übertrifft o1 zu 5\,\% der Kosten)
\end{itemize}

% ============================================================
% OPEN SOURCE (existiert)
% ============================================================
\section{Open Source -- Die Wettbewerber}

\subsection{Meta -- Llama-Familie}

Llama ist das dominierende Open-Source-Modell mit einem breiten Angebot an Größen:

\begin{table}[H]
\centering
\begin{tabularx}{\linewidth}{lXr}
\toprule
\textbf{Modell} & \textbf{Parameter} & \textbf{Use-Case} \\
\midrule
Llama 4 Maverick (1T) & 1 Billion (MoE) & Größte Qualität, Cloud-Inference \\
Llama 4 Scout (8B) & 8 Milliarden & On-Device, Edge, schnelle Response \\
Llama 3.3 (70B) & 70 Milliarden & Allrounder für eigene Infrastruktur \\
Llama 3.2 (1B/3B) & 1--3 Milliarden & Mobile und Embedded \\
\bottomrule
\end{tabularx}
\caption{Meta Llama-Modellfamilie -- Auswahl nach Einsatzzweck}
\label{tab:llama}
\end{table}

Llama ist unter einer modifizierten Apache-2.0-Lizenz verfügbar -- kommerzielle Nutzung kostenlos, bis zu einem bestimmten Umsatzlimit. Die 70B-Variante läuft gut auf einer einzelnen A100-GPU via vLLM.

\subsection{Mistral -- Das europäische Gegenstück}

Das französische Startup Mistral ist der führende europäische AI-Anbieter:

\begin{itemize}[leftmargin=*]
    \item \textbf{Mistral Large 2:} Das Flaggschiff, konkurrenzfähig mit Claude Sonnet. Besonders stark in europäischen Sprachen und Code-Generierung.
    \item \textbf{Mistral NeMo:} Mistral + NVIDIA Collaboration -- ein effizientes 12B-Modell mit 128K-Kontextfenster. Extrem kosteneffizient, open-weight.
    \item \textbf{Codestral:} Spezialisiert auf Code-Generation, besser als GPT-4o in Multi-Lang-Coding-Tasks.
\end{itemize}

Besonderheit: Mistral ist der einzige große europäische Anbieter mit globaler Reichweite und bietet \textbf{DSGVO-konforme Hosting-Optionen}.

\subsection{Qwen (Alibaba) und DeepSeek}

Chinesische Modelle gewinnen rasant an Qualität:

\begin{itemize}[leftmargin=*]
    \item \textbf{Qwen 2.5:} Besonders stark in Mathematik, Multilingualität und Code-Generation.
    \item \textbf{DeepSeek R1:} Der Preisbrecher -- Reasoning-Qualität auf o1-Niveau zu einem Bruchteil der Kosten.
    \item \textbf{DeepSeek V3:} Allgemeines Modell, konkurriert mit GPT-4o zu 10\,\% der Kosten.
\end{itemize}

% ============================================================
% CLOUD-PLATTFORMEN (existiert)
% ============================================================
\section{Cloud-Plattformen im Überblick}

Neben den nativen APIs bieten Cloud-Anbieter eigene KI-Plattformen:

\begin{table}[H]
\centering
\begin{tabularx}{\linewidth}{lX}
\toprule
\textbf{Plattform} & \textbf{Stärken} \\
\midrule
Azure AI (OpenAI) & Enterprise-Integration, bestehende Azure-Kunden, Compliance-Zertifikate \\
AWS Bedrock / SageMaker & Breiteste Modellauswahl, tiefste Cloud-Integration, eigene Inferenz \\
Google Vertex AI & Gemini-native, beste Multimodalität, GCP-Integration \\
Replicate & Einfaches Open-Source-Modell-Deployment (Serverless) \\
Together AI & Schnellstes Self-Hosting mit Llama \& Qwen \\
Groq & Extrem schnelle Inference (LPU-Technologie) für Open-Source-Modelle \\
\bottomrule
\end{tabularx}
\caption{Cloud-KI-Plattformen 2026}
\label{tab:cloud}
\end{table}

% ============================================================
% PRAXIS
% ============================================================
\section{Praxisbeispiele aus der Industrie}

\subsection{SaaS-Startup -- Dynamisches Modell-Routing}

Ein deutsches B2B-SaaS-Unternehmen mit 50.000 API-Calls pro Tag implementiert dynamisches Modell-Routing:

\begin{itemize}[leftmargin=*]
    \item Einfache Anfragen (Statusabfragen, Übersetzungen) $\rightarrow$ GPT-4o-mini
    \item Mittlere Anfragen (Analysen, Zusammenfassungen) $\rightarrow$ Claude Sonnet
    \item Komplexe Anfragen (Vertragsprüfung, Risikoanalyse) $\rightarrow$ GPT-4o
    \item Reasoning-Aufgaben (mathematische Optimierung) $\rightarrow$ DeepSeek R1
\end{itemize}

\textbf{Ergebnis:} 35\,\% Kostensenkung bei gleichbleibender Qualität. Fallback auf Anthropic bei OpenAI-Ausfällen.

\subsection{Automobilzulieferer -- Open Source aus Compliance-Gründen}

Ein deutscher Automobilzulieferer mit 20.000 Mitarbeitern setzt auf Llama 3 (70B) lokal gehostet:

\begin{itemize}[leftmargin=*]
    \item Grund: DSGVO und Geheimnisschutz. Keine Daten dürfen das Firmennetz verlassen.
    \item Betrieb: 4 A100 GPUs, vLLM, 200 Anfragen pro Minute.
    \item Ergebnis: Vergleichbare Qualität zu GPT-4o in der Domäne (technische Dokumentation).
    \item Kosten: Einmalige Anschaffung der Hardware vs. laufende API-Kosten.
\end{itemize}

\textbf{Lektion:} Ab etwa 1 Million Anfragen pro Monat ist Self-Hosting günstiger als API.

% ABBILDUNG: Kostenvergleich API vs Self-Hosted
\begin{verbatim}
Kosten (EUR/Monat)
^
|                          __ API-Kosten
|                        /
|                      /  ____ Self-Hosted (Fix)
|                    /  /
|                  /  /
|                /  /
|              /  /
|            /
|     _____/
+-------------------------------> Anfragen/Monat
      Break-Even ~ 1M
\end{verbatim}

\subsection{Fintech -- Multi-Provider als Risikoabsicherung}

Ein europäisches Fintech mit regulatorischen Anforderungen setzt auf drei Provider gleichzeitig:

\begin{itemize}[leftmargin=*]
    \item \textbf{Haupt-Provider:} Anthropic Claude (beste Quality-of-Service)
    \item \textbf{Backup:} OpenAI GPT-4o (bei Anthropic-Ausfällen)
    \item \textbf{Tertiary:} Selbst gehostetes Mistral (für sensible Daten)
    \item \textbf{Monitoring:} Latenz und Fehlerraten aller Provider in Echtzeit
\end{itemize}

\textbf{Ergebnis:} 99,95\,\% Verfügbarkeit des KI-Service über 12 Monate.

% ===========================================================%
% ENTSCHEIDUNGSMATRIX (existiert)
% ============================================================
\section{Entscheidungsmatrix: Welches Modell für welchen Fall?}

Diese Tabelle fasst die wichtigsten Entscheidungsgrundlagen zusammen:

\begin{table}[H]
\centering
\begin{tabularx}{\linewidth}{lXl}
\toprule
\textbf{Anforderung} & \textbf{Empfehlung} & \textbf{Alternativen} \\
\midrule
Beste allg. Qualität & Claude Sonnet 4 oder GPT-4o & Gemini Pro \\
Bestes Reasoning & o1-pro oder DeepSeek R1 & Claude Opus \\
Günstigste API & GPT-4o-mini oder Haiku & Qwen, Mistral NeMo \\
Längster Kontext & Gemini Ultra (1M Tokens) & GPT-4o (128K) \\
Bester Code & Codestral oder Claude Sonnet & DeepSeek-Coder \\
Datenhoheit (DSGVO) & Llama selbst gehostet & Mistral Enterprise \\
On-Device / Edge & Llama 3.2 1B/3B oder Gemini Nano & CoreML-Models \\
Multi-Modal (Bild+Audio) & GPT-4o oder Gemini Ultra & Claude (nur Text) \\
\bottomrule
\end{tabularx}
\caption{Modell-Empfehlungen nach Anforderungskategorie}
\label{tab:empfehlungen}
\end{table}

% ===========================================================%
% CODEBEISPIELE (existiert, erweitert)
% ============================================================
\section{Codebeispiele}

\subsection{Multi-Provider-Client für den Modell-Vergleich}

Ein praktischer AI Engineer wechselt oft zwischen Modellen -- daher ein Beispiel für eine abstrahierte API-Schicht:

\begin{lstlisting}[language=Python, caption={Abstrahierter Model-Client mit Multi-Provider-Support}, label=lst:multi-provider]
from abc import ABC, abstractmethod
from dataclasses import dataclass
from enum import Enum
import openai, anthropic


class Provider(Enum):
    OPENAI = openai
    ANTHROPIC = anthropic
    GOOGLE = google
    LOCAL = local

@dataclass
class QueryResult:
    response_text: str
    model_used: str
    tokens_input: int
    tokens_output: int
    cost_eur: float


class BaseModelClient(ABC):
    @abstractmethod
    def query(self, messages: list[dict]) -> QueryResult:
        pass

class OpenAIClient(BaseModelClient):
    def __init__(self, api_key: str, model: str = "gpt-4o-mini"):
        self.client = openai.OpenAI(api_key=api_key)
        self.model = model
        self.pricing = {"input": 0.0015, "output": 0.006}

    def query(self, messages: list[dict]) -> QueryResult:
        response = self.client.chat.completions.create(
            model=self.model,
            messages=messages,
            temperature=0.1,
        )
        return QueryResult(
            response_text=response.choices[0].message.content,
            model_used=self.model,
            tokens_input=response.usage.prompt_tokens,
            tokens_output=response.usage.completion_tokens,
            cost_eur=(response.usage.prompt_tokens * self.pricing["input"] +
                      response.usage.completion_tokens * self.pricing["output"]) / 1000,
        )

class AnthropicClient(BaseModelClient):
    def __init__(self, api_key: str, model: str = "claude-sonnet-4"):
        self.client = anthropic.Anthropic(api_key=api_key)
        self.model = model
        self.pricing = {"input": 0.003, "output": 0.015}

    def query(self, messages: list[dict]) -> QueryResult:
        response = self.client.messages.create(
            model=self.model, max_tokens=4096,
            messages=messages,
            system="Du bist ein hilfreicher Assistent."
        )
        return QueryResult(
            response_text=response.content[0].text,
            model_used=self.model,
            tokens_input=response.usage.input_tokens,
            tokens_output=response.usage.output_tokens,
            cost_eur=(response.usage.input_tokens * self.pricing["input"] +
                      response.usage.output_tokens * self.pricing["output"]) / 1000,
        )


def ask(messages: list[dict], provider: Provider = Provider.OPENAI):
    if provider == Provider.OPENAI:
        client = OpenAIClient(api_key="sk-...")
    elif provider == Provider.ANTHROPIC:
        client = AnthropicClient(api_key="sk-ant-...")
    result = client.query(messages)
    print(f"{result.model_used}: {result.response_text[:100]}... "
          f"(EUR{result.cost_eur:.4f})")

msgs = [{"role": "user",
         "content": "Was ist der Unterschied zwischen RAG und Fine-Tuning?"}]
ask(msgs, Provider.OPENAI)
ask(msgs, Provider.ANTHROPIC)
\end{lstlisting}

\subsection{TypeScript: Dynamischer Model-Router}

\begin{lstlisting}[language=TypeScript, caption={Dynamischer Model-Router in TypeScript}, label=lst:ts-router]
import OpenAI from "openai";
import Anthropic from "@anthropic-ai/sdk";

type Complexity = "simple" | "medium" | "complex" | "reasoning";

interface RouterConfig {
  simple: string;    // model name
  medium: string;
  complex: string;
  reasoning: string;
}

const config: RouterConfig = {
  simple: "gpt-4o-mini",
  medium: "claude-sonnet-4",
  complex: "gpt-4o",
  reasoning: "deepseek-r1",
};

async function routeRequest(
  prompt: string,
  complexity: Complexity,
): Promise<string> {
  const model = config[complexity];
  const client = new OpenAI();

  const resp = await client.chat.completions.create({
    model,
    messages: [{ role: "user", content: prompt }],
  });
  return resp.choices[0]?.message?.content ?? "";
}
\end{lstlisting}

% ===========================================================%
% BEST PRACTICES
% ============================================================
\section{Best Practices}

\begin{itemize}[leftmargin=*]
    \item \textbf{Abstrahiere den Model-Client}: Von Tag 1 an eine gemeinsame Schnittstelle für alle Provider. Der Wechsel kostet dann einen Konfigurations-Eintrag, kein Refactoring.
    \item \textbf{Führe ein Modell-Review durch}: Alle 3 Monate die aktuelle Modell-Landschaft prüfen. Preise ändern sich, neue Modelle kommen, alte verschwinden.
    \item \textbf{Nutze generische Modell-Namen nicht}: \texttt{gpt-4o} kann sich hinter den Kulissen ändern. Nutze datierte Versionen wie \texttt{gpt-4o-2026-03-15}.
    \item \textbf{Teste auf deinen Daten}: Ein Benchmark-Score ist kein Ersatz für einen Test mit deinen tatsächlichen Prompts und Erwartungen. Baue ein Golden Dataset.
    \item \textbf{Plane den Exit}: Jeder Provider muss austauschbar sein. Dokumentiere, welche Provider-spezifischen Features du nutzt (Tool-Use, Streaming, Batch-API).
\end{itemize}

% ===========================================================%
% ANTI-PATTERNS
% ===========================================================%
\section{Anti-Patterns}

\begin{itemize}[leftmargin=*]
    \item \textbf{Das beste Modell für alles}: Kein Modell ist universal am besten. Für 80\,\% der Tasks reicht GPT-4o-mini -- die anderen 20\,\% benötigen spezifische Stärken.
    \item \textbf{Vendor-Lock-in}: Zu tiefe Integration in einen Provider (z.\,B. Anthropic-spezifische Tool-Syntax) macht den Wechsel teuer. Abstrahierte Clients sind wertvolle Absicherung.
    \item \textbf{Vernachlässigung von Open Source}: Open-Source-Modelle sind 2026 so gut, dass sie für viele Use-Cases die API-Kosten komplett ersetzen können -- besonders mit Quantisierung.
    \item \textbf{Benchmark-Jagd}: Ein Modell nur nach MMLU-Score auswählen, ohne die eigenen Anforderungen zu testen.
    \item \textbf{Kein Fallback}: Ein einziger Provider, kein Backup. Fällt dieser aus, steht das gesamte System.
    \item \textbf{Immer das neueste Modell}: Neue Modelle haben oft Kinderkrankheiten. Warte 2--4 Wochen nach Release, bevor du in Produktion gehst.
\end{itemize}

% ===========================================================%
% PERFORMANCE
% ============================================================
\section{Performance und Kostenvergleich}

\begin{table}[H]
\centering
\begin{tabularx}{\linewidth}{lXXXX}
\toprule
\textbf{Modell} & \textbf{Input (pro Mio)} & \textbf{Output} & \textbf{Latenz (1K Tokens)} & \textbf{Max Kontext} \\
\midrule
GPT-4o-mini & 0,15 \$ & 0,60 \$ & 300--600 ms & 128K \\
GPT-4o & 2,50 \$ & 10,00 \$ & 800--2000 ms & 128K \\
Claude Sonnet 4 & 3,00 \$ & 15,00 \$ & 600--1500 ms & 200K \\
Claude Opus 4 & 15,00 \$ & 75,00 \$ & 1500--4000 ms & 200K \\
Gemini 2.5 Flash & 0,08 \$ & 0,30 \$ & 200--500 ms & 1M \\
Gemini 2.5 Pro & 1,25 \$ & 5,00 \$ & 500--1500 ms & 1M \\
DeepSeek R1 & 0,55 \$ & 2,19 \$ & 2000--8000 ms & 128K \\
Llama 3 (70B, lokal) & 0,00 \$ (Strom) & 0,00 \$ & 500--1500 ms & 128K \\
Mistral Large 2 & 2,00 \$ & 6,00 \$ & 600--1200 ms & 128K \\
\bottomrule
\end{tabularx}
\caption{Vollständiger Kosten- und Performance-Vergleich (Stand 2026)}
\label{tab:vollvergleich}
\end{table}

\textbf{Beobachtungen:}
\begin{itemize}[leftmargin=*]
    \item Der Kostenunterschied zwischen GPT-4o-mini und Claude Opus beträgt Faktor 100
    \item DeepSeek R1 bietet Reasoning auf o1-Niveau zu 10\,\% der Kosten
    \item Gemini Flash ist das günstigste API-Modell mit der niedrigsten Latenz
    \item Lokales Llama 3 amortisiert sich ab ca. 1 Mio. Anfragen/Monat
\end{itemize}

% ===========================================================%
% SICHERHEIT
% ===========================================================%
\section{Sicherheit -- Provisor-spezifische Risiken}

\subsection{Datenverarbeitung und Compliance}

\begin{table}[H]
\centering
\begin{tabularx}{\linewidth}{lX}
\toprule
\textbf{Anbieter} & \textbf{Datenverarbeitung} \\
\midrule
OpenAI & API-Daten standardmäßig 30 Tage gespeichert. Opt-out möglich (niedrigere Rate Limits). \\
Anthropic & API-Daten werden nicht zum Training verwendet. Standard bei Enterprise-Tier. \\
Google Gemini & Datenverarbeitung in GCP-Region wählbar. EU-Verarbeitung möglich. \\
Azure OpenAI & Vollständige EU-Datenhaltung. Vertragliche Auftragsverarbeitung. \\
Mistral & DSGVO-konformes Hosting in Frankreich/Deutschland. \\
Self-Hosted & Volle Kontrolle. Verantwortung für Sicherheit liegt beim Betreiber. \\
\bottomrule
\end{tabularx}
\caption{Datenverarbeitung und Compliance nach Anbieter}
\label{tab:compliance}
\end{table}

\subsection{Weitere Sicherheitsaspekte}

\begin{itemize}[leftmargin=*]
    \item \textbf{API-Key-Rotation}: Monatlicher Wechsel der API-Keys. Automatisierung über Secrets-Manager.
    \item \textbf{Rate-Limits}: Schützen vor versehentlicher Kostenexplosion (Loop ohne Limit).
    \item \textbf{Audit-Trail}: Jeder API-Call wird geloggt (User-ID, Prompt-Länge, Modell, Kosten).
    \item \textbf{Vendor-Risiko}: Was passiert, wenn der Provider insolvent geht oder sein Geschäftsmodell ändert? Exit-Strategie dokumentieren.
\end{itemize}

% ===========================================================%
% ZUSAMMENFASSUNG (existiert, erweitert)
% ============================================================
\section{Zusammenfassung}

\begin{itemize}[leftmargin=*]
    \item Die Closed-Source-Spitze besteht aus OpenAI (GPT-4o, o-Series), Anthropic (Claude) und Google (Gemini) -- jedes mit eigenen Stärken
    \item Llama (Meta), Mistral, Qwen und DeepSeek dominieren den Open-Source-Bereich und werden 2026 immer besser
    \item Cloud-Plattformen (Azure AI, AWS Bedrock, Vertex AI) bieten modellübergreifende Verwaltung und Enterprise-Features
    \item Die Modellwahl hängt stark vom Use-Case ab: Reasoning (o1), Qualität (Sonnet/4o), Kosten (mini/Haiku), Datenhoheit (Llama selbst gehostet)
    \item Ein Multi-Provider-Routing senkt Kosten um 30--50\,\% und erhöht Verfügbarkeit
    \item Vendor-Lock-in ist das größte Risiko -- abstrahierte Clients sind Pflicht
    \item Self-Hosting amortisiert sich ab ca. 1 Mio. Anfragen/Monat
\end{itemize}

% ===========================================================%
% MERKE-BOX
% ============================================================
\begin{merke}
\begin{itemize}
    \item \textbf{Kein Modell ist für alles optimal.} Die Kunst liegt im Routing.
    \item \textbf{GPT-4o-mini reicht für 80\,\% aller Aufgaben.} Teure Modelle nur bei konkretem Mehrwert.
    \item \textbf{Open Source ist 2026 gleichauf} mit Closed Source in den meisten Domänen.
    \item \textbf{Abstraktion ist Überleben.} Ein Multi-Provider-Client ist kein Luxus, sondern Versicherung.
    \item \textbf{Preise ändern sich monatlich.} Ein Review-Rhythmus von 3 Monaten ist Pflicht.
\end{itemize}
\end{merke}

% ===========================================================%
% PRAXISPROJEKT
% ============================================================
\section{Praxisprojekt}

\subsection{Aufgabe: Kostenoptimierung durch Modell-Routing}

Gegeben sind 1.000 Anfragen pro Tag mit folgender Verteilung:
\begin{itemize}[leftmargin=*]
    \item 70\,\% einfache Klassifikation (Modell A: GPT-4o-mini, 0,15 \$/Mio Input)
    \item 20\,\% mittlere Analyse (Modell B: Claude Sonnet, 3,00 \$/Mio Input)
    \item 10\,\% komplexes Reasoning (Modell C: DeepSeek R1, 0,55 \$/Mio Input)
\end{itemize}

Jede Anfrage verbraucht durchschnittlich 500 Input-Tokens und 200 Output-Tokens.

\textbf{Aufgaben:}
\begin{enumerate}[leftmargin=*]
    \item Berechne die monatlichen Kosten bei aktueller Verteilung
    \item Berechne die Kosten, wenn alle Anfragen über GPT-4o-mini laufen (Qualitätseinbußen ignoriert)
    \item Berechne die Kosten, wenn alle Anfragen über Claude Opus laufen (0\,\% Qualitätseinbußen)
    \item Wie viel sparst du durch das Routing im Vergleich zu \glqq nur Claude Opus\grqq{}?
\end{enumerate}

\textbf{Zusatz:} Baue ein Python-Skript, das diese Kosten dynamisch berechnet -- mit Konfiguration für Preise und Verteilung.

\textbf{Erfolgskriterium:} Du kannst die Kostenersparnis durch Routing beziffern und erklären, warum pauschale Modell-Entscheidungen teuer sind.

% ============================================================
% INTERVIEWFRAGEN
% ============================================================
\section{Interviewfragen}

\begin{enumerate}[leftmargin=*]
    \item \textbf{Ein Junior-Entwickler will ein \glqq 8x7B\grqq{} MoE-Modell (Mixture of Experts) auf einer kleinen 16-GB-GPU hosten. Sein Argument: \glqq Pro Token werden ja nur 14 Milliarden Parameter aktiviert, das passt locker.\grqq{} Wo liegt sein Denkfehler?}

    \textbf{Antwort:} Er verwechselt \textit{aktive} Parameter mit \textit{Speicherbedarf}. Damit das Modell routen kann, müssen \textit{alle} Gewichte des gesamten Modells im VRAM (Grafikspeicher) vorgehalten werden. Ein 8x7B-Modell benötigt (abzüglich geteilter Layer) rund 47 Milliarden Parameter. Selbst mit starker Quantisierung (INT4) sprengt das eine 16-GB-Karte bei weitem.

    \item \textbf{Du baust einen Semantic Router: Jede Anfrage geht zuerst an GPT-4o-mini zur Klassifikation, danach (wenn komplex) an Claude Opus. Das System funktioniert fehlerfrei, aber Nutzer beschweren sich über unerträgliche Wartezeiten. Was ist das architektonische Nadelöhr?}

    \textbf{Antwort:} Die sequenzielle Latenz. Die \glqq Time to First Token\grqq{} (TTFT) summiert sich. Wenn GPT-4o-mini 600~ms braucht und Opus danach nochmal 1500~ms TTFT hat, wartet der User über 2 Sekunden, bevor das erste Wort erscheint. Streaming löst das Problem nicht, wenn der Router den Stream blockiert. Bessere Lösung: Kleine lokale Klassifikatoren (FastText/BERT) mit 10~ms Latenz oder asynchrones Routing.

    \item \textbf{Ein neues Open-Source-Modell schlägt GPT-4o in den HumanEval- und MMLU-Benchmarks um Längen. In eurem System versagt es aber komplett daran, ein striktes JSON-Schema für eure API einzuhalten. Wie erklärst du dir das?}

    \textbf{Antwort:} Das ist ein klassischer Fall von \glqq Benchmark-Overfitting\grqq{} oder \glqq Data Contamination\grqq{}. Das Modell hat die Testdatenbanken während des Pre-Trainings auswendig gelernt, was die hohen Scores erklärt. Es mangelt ihm aber an echter \glqq Alignment\grqq{}-Qualität (Instruction Following) für komplexe Formatvorgaben wie JSON, die im Benchmark nicht in dieser Tiefe getestet wurden.

    \item \textbf{Euer System nutzt eine Fallback-Strategie: Wenn OpenAI einen \texttt{429 Rate Limit}-Fehler wirft, feuert das Backend den Request sofort exakt so gegen Anthropic. Nach einem viralen Marketing-Post gehen plötzlich beide Provider in eurer App offline. Warum?}

    \textbf{Antwort:} Es fehlt das \glqq Exponential Backoff \& Jitter\grqq{} sowie ein vernünftiges Queueing. Bei einem massiven Traffic-Spike laufen alle Requests gegen die Rate-Limits von Provider A und fluten im selben Millisekunden-Fenster Provider B, was dessen Limits sofort sprengt (Thundering Herd Problem). Solche Systeme erfordern Message-Broker (Kafka/RabbitMQ) zur Lastverteilung.

    \item \textbf{Der CFO will Kosten sparen und fordert, euren interaktiven Echtzeit-Support-Chatbot komplett auf DeepSeek R1 umzustellen, da es \glqq Reasoning auf o1-Niveau zu 10\,\% der Kosten\grqq{} bietet. Warum legst du Veto ein?}

    \textbf{Antwort:} Reasoning-Modelle (wie R1 oder o1) haben einen internen, unsichtbaren \glqq Chain-of-Thought\grqq{}-Prozess, bevor sie das finale Ergebnis ausgeben. Dadurch dauert es oft 5 bis 15 Sekunden, bis der erste Output-Token beim User ankommt. Das zerstört die UX in einem Echtzeit-Chat völlig. Solche Modelle sind für asynchrone Aufgaben (Agenten, Code-Generierung, komplexe Analyse) gedacht, nicht für den direkten Chat.

    \item \textbf{Euer Startup verarbeitet sensible Gesundheitsdaten (DSGVO). Du hast Mistral auf europäischen AWS-Servern vorgeschlagen, aber die Security-Abteilung lehnt ab: \glqq Die Daten verlassen unser VPC.\grqq{} Was ist die einzige echte Zero-Trust-Architektur?}

    \textbf{Antwort:} Das Self-Hosting eines Open-Weight-Modells (wie Llama 3 oder Qwen) auf komplett eigener, On-Premise-Hardware oder in einem isolierten Cloud-VPC (Virtual Private Cloud) \textit{ohne} Outbound-Internet-Zugang (Air-Gapped). Nur so ist technisch garantiert, dass keine Telemetrie- oder Log-Daten an Dritte abfließen.

    \item \textbf{Du abstrahierst OpenAI und Anthropic in einer gemeinsamen Client-Klasse. Ein Core-Feature deiner App benötigt zwingend den \texttt{json\_schema}-Modus von OpenAI für verschachtelte Daten. Anthropic bietet genau diesen Parameter nicht. Wie rettest du die Abstraktion?}

    \textbf{Antwort:} Ich mappe das Feature in meiner Middleware auf das \glqq Native Tool Calling\grqq{} (Function Calling) von Anthropic. Anthropic zwingt das Modell, Argumente für ein definiertes Tool (Funktion) zurückzugeben, was unter der Haube ein valides JSON nach vorgegebenem JSON-Schema erzwingt. So emuliere ich den Structured Output plattformunabhängig.

    \item \textbf{Ihr quantisiert ein 70B-Modell auf INT4, um es auf günstigeren GPUs laufen zu lassen. Im Smalltalk auf Deutsch funktioniert es super, aber es generiert plötzlich fehlerhaften Python-Code mit Syntax-Fehlern, was es in FP16 nicht tat. Warum?}

    \textbf{Antwort:} Quantisierung (das Kappen von Gleitkommazahlen-Präzision) trifft nicht alle Fähigkeiten gleich. Allgemeine Sprache ist sehr robust gegenüber Genauigkeitsverlusten. Komplexe Logik, Mathematik und Code-Syntax basieren jedoch auf hochsensiblen, feingetunten Gewichten. INT4 zerstört oft genau diese Randbereiche des Modellwissens.

    \item \textbf{Du tauschst das Backbone-Modell deines Agenten-Systems von GPT-4o auf Claude 3.5 Sonnet, verwendest aber exakt denselben 5.000 Tokens langen System-Prompt. Plötzlich weigert sich der Agent regelmäßig, völlig harmlose Aufgaben auszuführen. Warum?}

    \textbf{Antwort:} Modelle haben unterschiedliche Alignments (RLHF). Anthropics \glqq Constitutional AI\grqq{} macht Claude deutlich empfindlicher gegenüber potenziellen Safety-Verstößen (Refusal Rate). Zudem reagieren Modelle unterschiedlich auf Prompting-Stile (Claude bevorzugt XML-Tags zur Strukturierung, OpenAI Markdown). Ein System-Prompt ist Code und muss bei einem Modellwechsel immer modellspezifisch umgeschrieben und neu evaluiert werden.

    \item \textbf{Euer System fasst täglich 10.000 PDF-Verträge zusammen. Der Text ändert sich, aber der 2.000 Token lange System-Prompt mit den Extraktionsregeln bleibt immer gleich. Wie senkst du die Input-Kosten bei Anthropic oder OpenAI massiv, ohne das Modell zu wechseln?}

    \textbf{Antwort:} Durch \glqq Prompt Caching\grqq{}. Ich strukturiere den Prompt strikt so, dass die statischen 2.000 Tokens (Regeln, Few-Shot-Beispiele) immer exakt gleich am Anfang (als Prefix) stehen. Die APIs von Anthropic und OpenAI cachen diesen Prefix im KV-Cache, sodass für die System-Anweisungen nur noch ein Bruchteil (oft 10\,\% bis 50\,\%) der üblichen Input-Kosten berechnet wird.
\end{enumerate}
% ===========================================================%
% WEITERFÜHRENDE RESSOURCEN
% ===========================================================%
\section{Weiterführende Ressourcen}

\subsection{Papers und Analysen}

\begin{itemize}[leftmargin=*]
    \item \textbf{DeepSeek-R1: Incentivizing Reasoning Capability in LLMs via Reinforcement Learning} (DeepSeek, 2025)
    \item \textbf{The Llama 4 Herd of Models} (Meta, 2025)
    \item \textbf{Mixture-of-Experts Meets Instruction Tuning} (Mistral, 2024): Grundlagen zu MoE-Architekturen
\end{itemize}

\subsection{Dokumentationen}

\begin{itemize}[leftmargin=*]
    \item \href{https://platform.openai.com/docs}{OpenAI Platform Docs}
    \item \href{https://docs.anthropic.com}{Anthropic Documentation}
    \item \href{https://ai.google.dev}{Google AI Studio}
    \item \href{https://huggingface.co/models}{Hugging Face Model Hub}
    \item \href{https://ollama.ai/library}{Ollama Model Library}
\end{itemize}

\subsection{Open-Source-Projekte}

\begin{itemize}[leftmargin=*]
    \item \textbf{OpenRouter}: Einheitliche API für 100+ Modelle mit Fallback-Routing
    \item \textbf{LiteLLM}: Python SDK für 50+ Modelle mit einheitlicher Schnittstelle
    \item \textbf{vLLM}: Hochperformanter Inference-Server für Open-Source-Modelle
    \item \href{https://lmarena.ai}{Chatbot Arena}: Live-Benchmark mit menschlichen Bewertungen
\end{itemize}

\textbf{Nächster Schritt:} Im folgenden Kapitel vertiefen wir die harten Fakten: Kontextlänge, Knowledge Cutoffs und Pricing-Kalkulation -- die Faktoren, die im Alltag am meisten Kosten und Probleme verursachen.
